% Corrected French for printed p.149.
% Corrections: à caus de -> à cause de; celui de proposition -> celui de la proposition;
% Z^0 applied to the complex-valued resolving functor.
\SourcePage{149}

\LayerTarget{tohoku:section:2.5}
\noindent\textbf{2.5. Foncteurs résolvants.}

\LayerTarget{tohoku:def:resolving-functor}
Soient $C$ et $C'$ deux catégories abéliennes, on suppose que tout objet de
$C$ est isomorphe à un sous-truc d'un objet injectif. Soit $F$ un foncteur
covariant exact à gauche de $C$ dans $C'$. On appelle \textit{foncteur
résolvant} de $F$ un foncteur covariant $\mathbf F$ défini sur $C$, à valeurs
dans la catégorie des complexes à degrés positifs dans $C'$,
$\mathbf F(A)=(\mathbf F^n(A))$, muni d'une augmentation
$F\longrightarrow\mathbf F$ (i. e. d'un homomorphisme du foncteur $F$ dans
le foncteur $Z^0(\mathbf F)$ des cocycles en degré 0 de $\mathbf F$), et satisfaisant
aux conditions suivantes : (i) Le foncteur $\mathbf F$ est exact ;
(ii) $F\longrightarrow Z^0(\mathbf F)$ est un isomorphisme ; (iii) si $A$ est
injectif, alors $\mathbf F(A)$ est acyclique en degrés $>0$.

Soit $\mathbf F$ un foncteur résolvant pour $F$. Considérons les foncteurs
$H^n\mathbf F(A)$ sur $C$, à valeurs dans $C'$. À cause de la condition (i),
ils forment un foncteur cohomologique, à cause de (ii) ce dernier se réduit à
$F$ en dimension 0, enfin en vertu de (iii) les $H^n\mathbf F(A)$ pour $n>0$
sont effaçables. Par suite :

\LayerTarget{tohoku:prop:2.5.1}
\noindent\textsc{Proposition 2.5.1.} --- \textit{Si $\mathbf F$ est un
foncteur résolvant pour $F$, alors on a pour tout $A\in C$ des isomorphismes
uniques $H^n\mathbf F(A)=R^nF(A)$, définissant un isomorphisme de foncteurs
cohomologiques, et se réduisant en dimension 0 à l'isomorphisme
d'augmentation.}

Nous allons donner une autre démonstration de la proposition précédente,
permettant un calcul commode de ces isomorphismes :

\LayerTarget{tohoku:prop:2.5.2}
\noindent\textsc{Proposition 2.5.2.} --- \textit{Soit $\mathbf F$ un foncteur
résolvant pour $F$. Soit $A\in C$, et soit $C=(C^p(A))$ une résolution droite
de $A$ par des objets $F$-acycliques. Considérons le bicomplexe
$\mathbf F C(A)=(\mathbf F^qC^p(A))_{p,q}$, et les homomorphismes naturels}
\LayerTarget{tohoku:eq:2.5.1}
\[
\tag{2.5.1}
F(C(A))\longrightarrow\mathbf F(C(A))\longleftarrow\mathbf F(A)
\]
\textit{définis respectivement par l'augmentation
$F\longrightarrow\mathbf F$ et par l'augmentation
$A\longrightarrow C(A)$. Alors les homomorphismes correspondants}
\LayerTarget{tohoku:eq:2.5.2}
\[
\tag{2.5.2}
H^nF(C(A))=R^nF(A)\longrightarrow H^n\mathbf F(C(A))
\longleftarrow H^n\mathbf F(A)
\]
\textit{sont des isomorphismes, et l'isomorphisme correspondant
$H^n\mathbf F(A)=R^nF(A)$ est celui de la proposition 2.5.1.}

\LayerTarget{tohoku:prop:2.5.2:proof:p1}
Nous considérons $F(C(A))$ (resp. $\mathbf F(A)$) comme un bicomplexe dont le
deuxième (resp. premier) degré est nul. La première suite spectrale de
$\mathbf F(C(A))$ a pour terme initial $H(F(C(A)))$ (car
$H_{\mathrm{II}}^q(\mathbf F(C(A)))$ est nul pour $q>0$ et s'identifie à
$F(C(A))$ pour $q=0$, en vertu de prop. 2.5.1 et du fait que les $C^p(A)$ sont
$F$-acycliques), donc le premier homomorphisme (2.5.1) induit un isomorphisme
pour les termes initiaux de la suite spectrale $I$, donc induit un
isomorphisme pour la cohomologie. De même,
$H_{\mathrm I}(\mathbf F(C(A)))^{p,q}$ est nul pour $p>0$ et égal à
$\mathbf F(A)$ pour $p=0$ (puisque $\mathbf F$ est un foncteur exact et
$C(A)$ est une résolution de $A$), donc le deuxième homomorphisme (2.5.1)
induit un isomorphisme pour les termes initiaux des suites spectrales $II$,
donc induit encore un isomorphisme pour la cohomologie. Pour montrer que
l'isomorphisme obtenu de $H^n\mathbf F(A)$ sur $R^nF(A)$ est bien celui de la
proposition 2.5.1, on peut prendre une résolution \textit{injective} $C'(A)$
de $A$, un homomorphisme de $C(A)$ dans $C'(A)$, et envisager
l'homomorphisme associé du diagramme (2.5.2)

\TohokuContinuation{150}
